\section{Numbers According to Iamblichus}

Plato said that no one could be a philosopher who had not studied mathematics.

Undoubtedly, this is partially a reference to the Pythagoreans. Having been a mathematical dunce until college (although not innumerate entirely), this is part of my penance, to work through Iamblichus' treatise.

Luckily, this work supplies a beginner's manual to understanding the significance of number, which plays a role in traditional astrology. Still, there is a lot to wade through here. For instance, John Michael Greer (the so-called Archdruid of North America) has \emph{pointed} out \emph{that the Christian cross\footnote{\url{https://en.wikipedia.org/wiki/Sator\_Square}} has the proportions of the square root of 2\footnote{\url{http://www.geomancy.org/sacred-geometry/square-root-3/index.php\#sacred-geometry/square-root-2/index.php}} and the square root of 5\footnote{\url{www.geomancy.org/sacred-geometry/square-root-3/index.php\#sacred-geometry/square-root-2/index.php}}.} The\footnote{\url{http://www.keyofsolomon.org/True\%20cross.php}} cross is an ancient symbol\footnote{\url{http://www.ccg.org/english/s/p039.html}}, one which Plato was familiar with when he wrote \emph{``the world soul is crucified''.} However, it is not necessary (obviously) to be a Christian to discern the Logos inherent in the numbers. In our day, it may help if one is not exoterically Christian to see the significance of the numbers in constituting a pattern of the Logos.

Iamblichus touches on a great deal of this, in how the numbers intertwine, which require a beginner level in numeracy, but patience in working through the various relationships, and a meditative or mystical approach to thinking about them.

\begin{enumerate}
\item Any number can be created by adding Ones (even as fractions) or dividing by Ones, or multiplying by Ones. The One is therefore the Source and Sustainer of all things. If x = 2y, then x squared is four times y squared, which proportion in operation is preserved if one take 1/x and 1/y. There is no change in the nature or operations of the One. It is called the Provider because it has the power of staying the same, regardless of its extensions, and preserving both itself and those extensions. The ancients called it Prometheus or the Artificer of Life.

\item It is called the Hearth, because it is in the middle of things, that which is equal between opposing elements, and it keeps its equilibrium. In the middle of the four elements lies a Monadic fiery cube. It is called Proteus, because it assumes any form, but retains the properties of everything, as the Monad is the factor of each number.

\item It is called Chaos (mixture, obscurity, blending, Darkness), being Hesiod's First Generator, but also called Chariot, Sun, Friend, Life, Happiness, Being, Order, and Concord. The numerical value of the Greek word Monad adds up to 361, the degrees of the circle, plus One. It is disposed to share itself with all things. It is not yet manifesting anything actual, but carries within it the principles conceptual togetherness of all things. It causes things to combine and co-here, and unifies that which is opposite. It is called Androgyne and Intellect, as well, and shares power with the Sun. It is the beginning, middle, and end of all size, quality and quantity.

\item It produces itself out of itself, as well as producing all things, and is therefore ``as if it were God or the principle of all things''. It maintains everything and forbids anything to change --- it resembles Providence, alone of all numbers. It is particularly suited to resemble and reflect God, and it is closest to Him.

\item It is in fact the Form of Forms. As we will see, the Dyad is opposed to Form in a certain sense, and only from the One can the power come to maintain the Dyad. Because of this intellective and creative power, it is both Creator and Supreme Intellect.
\end{enumerate}

What else is Iamblichus describing here, but a doctrine of the Logos? This Logos would be taken up and shaped by Christianity, and I would argue that the intention was to preserve ``every jot and tittle'' of the full-blooded pagan dogma of the Logos, rather than to change it in any respect, except to allow it to be ``more itself''.

And it is clear that he thinks that multiplication by One preserves the same number is a truth with an esoteric meaning. Or, division. Mathematics are not, therefore, just mathematics, even the simple operations we don't even think about. Yet Iamblichus will get into more complicated arrangements as he progresses.


\flrightit{Posted on 2012-12-28 by Logres}

